\begin{dang}{Phương trình Cla-pê-rôn - Mendeleev (dành cho ban nâng cao)}
	\Bookmark Phương trình CLA-PÊ-RÔN - MEN-ĐÊ-LÊ-ÉP:
	\begin{align*}
		\boder{p V = n RT =\dfrac{m}{\mu}RT}\quad
		\ct{Trong đó: }
		\begin{cases}m: \ct{khối lượng chất khí (g)}\\
			\mu: \ct{khối lượng mol của chất khí (g/mol)}\\
			n: \ct{số mol chất khí}\\
			T: \ct{nhiệt độ (K)}\\
			P: \ct{áp suất (Pa hoặc $N/m^2$)}\\
			V: \ct{thể tích ($m^3$)}\\
			R: \ct{hằng số khí ($R = 8,31$ J/mol.K)}
		\end{cases}
	\end{align*}
	\textbf{Chú ý:}
	\begin{itemize}
		\item Trong trường hợp $p$ là atm và V là lít thì $R = 0,082\ atm.\ell/mol.K$.
		\item Trong trường hợp p là Pa hoặc N/m thì $R = 8,31\ J/mol.K$.
		\item Lực đẩy Acsimet của không khí: 
		\begin{align*}F =  DVg= \rho V\quad  \ct{Với} \begin{cases}\ct{D là khối lượng riêng }(kg/m^3)\\ \rho \ct{ là trọng lượng riêng } (N/m^3)\\ \ct{V là thể tích mà vật chiếm chỗ}.
			\end{cases}
		\end{align*}
	\end{itemize}
\end{dang}